\section{Nixie tube characteristics}
\label{sec:tubes}

The design carries six cold-cathode gas-discharge indicator tubes, forming
an HH:MM digit display and a colon separator between the hour and minute
pairs. This document assigns the designators used throughout to refer to
them:

\begin{table}[htbp]
\centering
\caption{Nixie tube designators}\label{tab:tube-designators}
\begin{tabular}{|l|l|}
\hline
\textbf{Designator} & \textbf{Component} \\
\hline
NX1--NX4 & IN-12B digit tubes (HH:MM display) \\
\hline
NX5--NX6 & INS-1 colon tubes (colon separator) \\
\hline
\end{tabular}
\end{table}

Both are Soviet-made parts documented only in scanned factory sheets;
the figures below are transcribed directly from those scans, with the
original Russian parameter names noted where a term's meaning is not
self-evident from context. Neither sheet carries guaranteed
min/typ/max columns throughout. Several entries below are single-sided
limits (``not more than'' or ``not less than''), which bears on how
conservatively the anode supply must be sized.

\subsection{IN-12B (digit tubes, NX1--NX4)}
\label{sec:in12b}

Source: \texttt{IN-12A\_IN-12B.pdf}, referenced below as [RU] (original
Soviet factory sheet, covering both IN-12A and IN-12B; IN-12B adds a
decimal-point cathode that IN-12A lacks), and
\texttt{IN-12\_english\_gra-afch.pdf}, referenced below as [EN] (an
English-language summary sheet, used as a cross-check).

{\footnotesize
\begin{longtable}{|L{3.6cm}|c|c|c|c|L{4.4cm}|c|}
\caption{IN-12B electrical and physical characteristics}\label{tab:in12b}\\
\hline
\textbf{Parameter} & \textbf{Min} & \textbf{Typ} & \textbf{Max} & \textbf{Units} & \textbf{Conditions} & \textbf{Source} \\
\hline
\endfirsthead
\caption[]{IN-12B electrical and physical characteristics (continued)}\\
\hline
\textbf{Parameter} & \textbf{Min} & \textbf{Typ} & \textbf{Max} & \textbf{Units} & \textbf{Conditions} & \textbf{Source} \\
\hline
\endhead
Maintaining (plate) voltage & & \num{145} & & \si{\volt} & ``Typical Plate Voltage'' & [EN] p.1 \\
\hline
Ignition voltage & & & \num{170} & \si{\volt} & discharge-onset voltage (``Напряжение возникновения разряда''); not split by variant & [RU] p.4 \\
\hline
Cathode (digit) current, DC operating & \num{2.5} & & \num{3.5} & \si{\milli\ampere} & ``Ток рабочий, постоянный, для цифр'' & [RU] p.4 \\
\hline
Cathode (digit) current, average (unfiltered half-wave AC) & \num{1} & & \num{2} & \si{\milli\ampere} & ``Ток рабочий, средний\ldots для цифр'' & [RU] p.4 \\
\hline
Cathode (digit) overload current & & & \num{5} & \si{\milli\ampere} & Time-limited; see overload duration below & [RU] p.4 \\
\hline
Decimal-point cathode current, DC operating & \num{0.3} & & \num{0.5} & \si{\milli\ampere} & IN-12B only (pin 12); IN-12A pin 12 unconnected & [RU] p.4 \\
\hline
Decimal-point overload current & & & \num{0.7} & \si{\milli\ampere} & & [RU] p.4 \\
\hline
Overload duration limit, IN-12A & & & \num{30} & min & Applies to the \SI{5}{\milli\ampere} digit overload above & [RU] p.5 \\
\hline
Overload duration limit, IN-12B & & & \num{33} & min & Applies to the \SI{5}{\milli\ampere} digit / \SI{0.7}{\milli\ampere} decimal-point overload above & [RU] p.5 \\
\hline
Anode supply voltage & \num{200} & \num{250} & & \si{\volt} & DC or half-wave-rectified AC; the \num{250}~V recommended value reduces discharge-onset delay time & [RU] p.4, p.5 \\
\hline
Conditioning current, pulsed & \num{1} & & \num{2} & \si{\milli\ampere} & $\geq$\SI{1}{\minute}/cathode, applied once after storage & [RU] p.5 \\
\hline
Conditioning current, DC & \num{2.5} & & \num{3} & \si{\milli\ampere} & $\geq$\SI{1}{\minute}/cathode, applied once after storage & [RU] p.5 \\
\hline
Continuous single-cathode life rating & & & \num{1000} & h & ``Longevity'' mode, one cathode continuously lit & [RU] p.5 \\
\hline
Digit height & & \num{18} & & mm & & [EN] p.1 \\
\hline
Envelope diameter & \num{29} & & \num{31} & mm & $-2$~mm tolerance [RU]; [EN] gives \SI{32}{\milli\meter} & [RU] p.3, [EN] p.1 \\
\hline
\end{longtable}
}

\paragraph{Not specified in either source.} A combined (total) anode
current limit distinct from the per-cathode figures; an interelectrode
or anode-cathode capacitance; a manufacturer-recommended ballast
resistor value or target anode current within the permissible 2.5--3.5~mA
band.

\paragraph{Open question: anode supply voltage.} The datasheet's
minimum anode supply figure (\SI{200}{\volt}, recommended
\SI{250}{\volt}) sits \emph{above} the \SIrange{164.8}{191.5}{\volt}
anode-rail range this document derives (\cref{sec:hv-converter}), so
the design does not reach even the stated minimum. The likely
explanation is that the factory figure is written for the tube's other
stated supply mode, unfiltered half-wave-rectified AC
(\S\ref{sec:in12b}, ``Cathode current, average''). An unsmoothed
half-wave supply spends roughly half of every mains cycle near zero
volts, so its \emph{peak} must be pushed well above the
\SI{145}{\volt} typical maintaining voltage to sustain adequate
\emph{average} current and reliable strike. A continuously present,
regulated DC rail carries no such requirement. The source does not
explain the discrepancy, so this reading is unconfirmed.
\Cref{sec:hv-converter} carries it forward as an open item; the
bench-tuning procedure (raising the trim potentiometer until every
digit of the actual tube stock strikes reliably, then setting the rail
roughly \SI{10}{\volt} above that point) tests the uncertainty
directly, whichever reading is correct.

\subsection{INS-1 (colon tubes, NX5--NX6)}
\label{sec:ins1}

Source: \texttt{INS-1.pdf} (original Soviet factory label sheet,
ODO.334.095 TU). The tube has two electrode leads, counted from the
indicator dot: lead 1 is the anode (A, the small cylinder) and lead 2
the cathode (K, the large cylinder). This assignment follows the
datasheet's electrode table. There is no auxiliary or priming
electrode.

{\footnotesize
\begin{longtable}{|L{4.6cm}|c|c|c|c|L{3.6cm}|c|}
\caption{INS-1 electrical characteristics}\label{tab:ins1}\\
\hline
\textbf{Parameter} & \textbf{Min} & \textbf{Typ} & \textbf{Max} & \textbf{Units} & \textbf{Conditions} & \textbf{Page} \\
\hline
\endfirsthead
\caption[]{INS-1 electrical characteristics (continued)}\\
\hline
\textbf{Parameter} & \textbf{Min} & \textbf{Typ} & \textbf{Max} & \textbf{Units} & \textbf{Conditions} & \textbf{Page} \\
\hline
\endhead
Striking voltage & \num{65} & & \num{90} & \si{\volt} & Discharge-onset voltage (``Напряжение возникновения разряда''); no test current given & 2 \\
\hline
Maintaining voltage & & & \num{55} & \si{\volt} & ``Напряжение поддержания разряда''; single-sided limit, no minimum given & 2 \\
\hline
Optimal indication current & & \num{0.5} & & \si{\milli\ampere} & ``Оптимальным током, обеспечивающим достаточную яркость свечения''; a brightness recommendation, with no band around it & 3 \\
\hline
\end{longtable}
}

\paragraph{Conditioning after storage.} The datasheet directs that an
indicator taken out of long storage be run at its nominal discharge
current for one to two minutes before normal service, and refers
operation to GOST 11163-84 and OST 11~339.003-75. The IN-12B carries a
comparable instruction (\cref{sec:in12b}).

\paragraph{Not specified.} The sheet gives no minimum for the
maintaining voltage, no rated operating-current band around the
\SI{0.5}{\milli\ampere} brightness figure, no maximum discharge
switching frequency, no interelectrode capacitance, no example ballast
value, and no envelope dimensions. Two consequences follow for the
derivations downstream. \Cref{sec:hv-load-colon} sizes the ballast
against the single \SI{0.5}{\milli\ampere} figure and reports the
spread the trim range produces around it, having no band to check
against. And because \SI{55}{\volt} is a ceiling with no floor beneath
it, a tube maintaining below that figure draws proportionally more
current through a fixed ballast, which that section treats as the open
side of the result.
